%% kapa: exam=1 round=25 year=2014 period=2 form=A subject=appraisal_law subject_ko=부동산관계법규 seq=21 number=21
%% source: Q-Net 감정평가사 기출문제 (cst00309) artlSeq=5101207 file=extracted/1008/5101207/2014년 제25회 감정평가사 1차시험 기출문제/2교시(A형).hwp
%% conversion: kordoc 4.24.51 -> tools/gen1.py (auto); review=proofread-sonnet; flags=none
\begin{question}{21}
부동산 가격공시 및 감정평가에 관한 법령상 감정평가법인에 관한 설명으로 옳은 것은?
\begin{choices}
  \item 감정평가사가 아닌 자는 감정평가법인의 대표이사가 될 수 없다.
  \item 감정평가법인에는 15인 이상의 감정평가사를 두어야 한다.
  \item 감정평가법인이 주주총회의 의결에 따라 다른 감정평가법인과 합병하는 경우에는 합병 후 7일 이내에 국토교통부장관에게 이를 신고하여야 한다.
  \item 감정평가법인의 주사무소에 주재하는 법정 최소 감정평가사의 수는 2명이다.
  \item 감정평가법인의 자본금은 3억원 이상이어야 한다.
\end{choices}
\end{question}
\answer{4}
